\documentclass[11pt,a4paper]{article}
\usepackage[margin=25mm]{geometry}
\usepackage{amsmath,amssymb,amsthm,booktabs,longtable,array,hyperref}
\newtheorem{theorem}{Theorem}
\newtheorem{proposition}[theorem]{Proposition}
\newtheorem{lemma}[theorem]{Lemma}
\newcommand{\gen}[1]{\langle #1\rangle}
\newcommand{\Aut}{\operatorname{Aut}}
\newcommand{\Z}{\mathbb Z}
\title{Locally spherical regular hypertopes:\\marked quotients, affine families and a corrected finite table}
\author{Human Compression project}
\date{}
\begin{document}\maketitle
\begin{abstract}
We give a complete attributed account of the fixed revised manuscript
\emph{An exploration of locally spherical regular hypertopes} by Fernandes,
Leemans and Ivi\'c Weiss. The spherical classification follows from Maxwell's
retained normal-subgroup theorem and explicit antipodal incidence arguments.
A reflection-commutator proof gives the affine translation-kernel conclusion
under the universal-residue hypothesis, and the displayed toroidal families
are joined to their exact parameter domains. A finite separating-set argument
supplies arbitrarily large finite examples for every retained compact
hyperbolic diagram of rank at least three. Complete bounded quotient
certificates verify twenty-two named rows of the source's finite table and
exclude its other four rows. We state both diagram-only and universal-residue
minimum conventions. Auxiliary quotient assertions are corrected explicitly;
classifications and earlier constructions retained from the literature are
identified throughout.
\end{abstract}

\section{Scope, conventions and acquired proofs}\label{scope}
Our source is the sixteen-page revised manuscript~\cite{FLW}, whose fixed
PDF identity and complete clause ledger accompany this article. Its page~5
defines a spherical hypertope by a finite Coxeter diagram; this allows
projective quotients. Section~5, page~10, expressly restricts instead to
universal spherical residues. We call the first condition \emph{broad} and
the second \emph{locally universal}. Neither convention is silently imposed
on the source's hyperbolic table. Isomorphism in the source permits reordering
types, including duality. Counts of normal kernels for an ordered presentation
are a different invariant and will never be called counts of unlabelled
hypertopes.

A hypertope is a thin, residually connected incidence geometry. Its group
$\Aut_I$ of type-preserving automorphisms is distinguished from its group
allowing type permutations. Regularity means that $\Aut_I$ acts regularly
on chambers. For an involutory marking $G=\gen{r_i:i\in I}$ write
$G_J=\gen{r_j:j\in J}$ and $P_i=G_{I\setminus\{i\}}$. The intersection
property is $G_J\cap G_K=G_{J\cap K}$ for all $J,K\subseteq I$.

The complete arguments acquired in this article are supplied as separate
editable manuscripts in the reading volume, with their finite certificates
and distinct reviews. They are: the residue and marking theorem~\cite{residue};
the quotient-kernel correction~\cite{kernel}; the spherical argument~\cite{spherical};
the affine kernel and period arguments~\cite{translation,period}; finite
residue-preserving quotients~\cite{existence}; the eight-diagram finite core
and named groups~\cite{core,names}; the exact Coxeter/LINS interfaces~\cite{methods};
the ninth-square and exact-order exclusions~\cite{square,divisor}; and the
infinite cubic example~\cite{cubes}. The present proofs make the remaining
elementary source joins. Acquiring these complete proofs does not reprove
their explicitly stated classical or computational premises.

\section{Markings, cosets and the quotient correction}\label{markings}
We use the proved residue theorem~\cite{residue}: the distinguished chamber
involutions of a regular hypertope generate its full type-preserving group,
satisfy the intersection property, and identify all flag stabilizers with
the corresponding marked subgroups. Conversely a C-group whose actual coset
incidence system is flag transitive gives a regular hypertope. This is the
interface of source Theorem~2.3; the string C-group/polytope correspondence
in source Theorem~2.4 remains the classical Aschbacher--Schulte result.

\begin{proposition}\label{cosets}
For any family $(P_i)_{i\in I}$ of subgroups of a group $G$, the typed
right cosets $P_i g$, with incidence given by nonempty intersection, form
an incidence system with base chamber $(P_i)_{i\in I}$. Right multiplication
acts transitively on flags of each prescribed type of size at most two.
Its kernel is $\bigcap_i\operatorname{Core}_G(P_i)$.
\end{proposition}
\begin{proof}
Cosets of the same subgroup are equal or disjoint, which verifies the
same-type incidence axiom. All base cosets contain the identity.
Multiplication preserves intersections. For an incident pair choose an
element $g$ in its intersection; the pair is then $(P_i g,P_j g)$.
This also proves transitivity for singleton flags. An element fixes every
coset of type $i$ precisely when it lies in every conjugate of $P_i$.
Intersecting this condition over the types gives the asserted kernel.
\end{proof}
Thus the sentence preceding source Proposition~2.2 must not identify an
arbitrary $G$ with the full automorphism group. The proposition's actual
action statement is valid. For a regular hypertope the stabilizer/coset
theorem supplies the faithful full group, not this general construction alone.

Source Proposition~2.1 and Theorem~3.1 are proved in~\cite{kernel}.
In particular, C-group maximal parabolics and all the expected pairwise
maximal intersections imply the full intersection property. An epimorphism
to a C-group that is injective on all maximal marked parabolics therefore
lifts that property. This implication points \emph{upward}; it does not
give a reverse homomorphism from a quotient to its cover.

The exact downward criterion, under local injectivity, is
\begin{equation}\label{pair}
 N\cap P_i=1\quad\hbox{for all }i,\qquad
 G/N\hbox{ is a C-group} \Longleftrightarrow
 N\cap P_iP_j=\{1\}\quad(i\ne j).
\end{equation}
Here $G$ itself is a C-group and $N\triangleleft G$; the proof lifts two
equal coset images into $P_i$ and $P_j$ and applies the maximal-parabolic
criterion. Flag transitivity must still be checked. The affine $\widetilde A_2$
example in~\cite{kernel} has $G=\Z^2\rtimes S_3$ and $N=\Z^2$.
Each maximal parabolic maps injectively onto $S_3$, yet the quotient
incidence system has one object of each type and one chamber. It is flag
transitive but not thin. This is a counterexample to source Theorem~3.2,
and~\eqref{pair} is its precise intersection repair.

For any regular hypertope, its marked Coxeter cover $W\twoheadrightarrow G$
and coset reconstruction identify the hypertope with the incidence quotient
of the universal complex. Local injectivity follows if its proper residues
are universal, since their chamber groups have the same finite orders as
the covering parabolics. It does not follow from a finite residue diagram
alone. The order-$1920$ square in Section~\ref{finite} has local groups of
order $60$, whereas the covering $H_3$ parabolics have order $120$.
It is an explicit broad counterexample to the local-injectivity clause of
source Theorem~3.4. The marked quotient and incidence reconstruction survive.

\begin{proposition}\label{reducible}
A regular hypertope whose diagram has connected components $I_1,\ldots,I_t$
is the direct sum of the corresponding regular hypertopes. Its distinguished
group is $G_{I_1}\times\cdots\times G_{I_t}$.
\end{proposition}
\begin{proof}
Marks in different components commute. The intersection property gives
$G_{I_a}\cap G_{I\setminus I_a}=1$. Multiplying the commuting factors
therefore gives an injective homomorphism from their direct product onto
$G$. For $i\in I_a$ the maximal parabolic equals the product of
$(G_{I_a})_{I_a\setminus\{i\}}$ and all other factors. Its cosets are
thus exactly the cosets in the $a$th factor. Cosets belonging to different
components always intersect; within a component their incidence is its
factor coset incidence. This proves the direct-sum description. Fixing the
base objects outside $I_a$ identifies that factor with a residue of the
original hypertope. It is consequently a regular hypertope by the acquired
residue theorem.
\end{proof}

\subsection*{The two geometric examples}
For source Example~2.5, colour unit cubes of $\Z^3$ by the parity of their
lower vertex. The exact dual-cell construction~\cite{cubes} sends the two
cube colours to even and odd vertices of the shifted dual grid, source
edges to dual squares, and source vertices to dual cubes. Source types
$(0,2,1,3)$ become types $(0,1,2,3)$. The direct infinite-grid proof gives
flag extension, thinness, residual connectivity and a regular chamber
action of $D\rtimes B_3$, where
$D=\{v\in\Z^3:\sum v_i\equiv0\pmod2\}$. Its marks are
$\tau_{(1,1,0)}z_1z_2s_1,s_1,s_2,z_3$, with $z_i$ sign changes and
$s_i$ adjacent coordinate swaps. This is the zero-period object itself;
it is not obtained by setting a finite parameter restricted to $t\ge2$ to zero.

For precision, the negative flag in~\cite{spherical} also proves the
nongeometry assertion of source Example~3.3. In the usual $D_4$ model
let $v=(1,1,1,1)/2$, $w=(1,1,1,-1)/2$, and $x=e_4$. Their antipodal
classes are pairwise incident. A fourth-type object is represented by
$y=\pm e_i\pm e_j$. Incidence with $[x]$ requires that one coordinate
be $4$. Incidence with $[v]$ then requires the two signs of $y$ to agree,
whereas incidence with $[w]$ requires them to disagree. Thus this flag
has no fourth-type extension. This strengthens the same signed-triangle
obstruction from failure of the proposed regular action to failure of
flag extension in this rank-four example.

\section{The complete spherical classification}\label{sphericalsec}
Retain Maxwell's Theorem~0.1~\cite{Maxwell}, as quoted in source
Theorem~4.1: a normal subgroup of an irreducible finite Coxeter group is
the central antipodal subgroup or the kernel of a graph homomorphism.
The primary full proof is not acquired. A graph homomorphism sends simple
generators to simple generators or to the identity and covers the target
simple generators. In a hypertope quotient the marks remain distinct,
nontrivial, and preserve every pair order. Such a graph homomorphism is
an isomorphism of the defining Coxeter presentations. The only remaining
kernel is the central antipode, when present.

The full argument~\cite{spherical} proves that the antipode avoids every
proper parabolic and every product $P_iP_j$. It therefore gives a C-group
quotient. The string correspondence pays its geometry for string diagrams.
For the other cases, fundamental-weight inner products determine incidence;
every quotient flag lifts exactly when its signed triangles have positive
sign product. A single explicit triangle excludes every even $D_n$ case.
For $E_7,E_8$, fifteen complete, generator-closed weight orbits, containing
$899402$ vectors, exclude every bad triangle. Their accepted certificates
are acquired, not enumerated again here. The resulting classification is:
\begin{center}\small
\begin{tabular}{llrc}\toprule
Diagram & Range & Universal group order & Additional antipodal object\\\midrule
$A_n$ &$n\ge1$&$(n+1)!$&no\\
$D_n$ &$n\ge4$&$2^{n-1}n!$&no\\
$B_n$ &$n\ge3$&$2^n n!$&yes\\
$I_2(p)$ &$p\ge3$&$2p$&no\\
$E_6$ &&$51840$&no\\
$E_7$ &&$2903040$&yes\\
$E_8$ &&$696729600$&yes\\
$F_4$ &&$1152$&yes\\
$H_3$ &&$120$&yes\\
$H_4$ &&$14400$&yes\\\bottomrule
\end{tabular}
\end{center}
An additional antipodal object has half the displayed group order. All
proper residues of every object in this table are universal spherical,
by the embedded parabolics and exact residue reconstruction. This proves
source Theorems~4.2--4.3, subject to the named retained inputs, without
assuming local injectivity in order to obtain it.

The $E_6$ entry corrects source Table~1's $12\cdot6!$ to $72\cdot6!$.
Its exact certificate has a $27$-element fundamental-weight orbit with
standard $D_5$ stabilizer of order $1920$, giving $27\cdot1920=51840$.
Conventional extended Schl\"afli notation in source Table~1 is retained
as notation for these universal/projective objects. No additional
presentation theorem is inferred from that notation. The source's blanket
half-power description of a longest element is not used: in standard
$A_3$, $c=(1234)$ has $c^2=(13)(24)$, while the longest element is
$(14)(23)$. Only actual central antipodes, with the exact powers verified
in~\cite{spherical}, occur in the proof.

\section{Affine kernels and the displayed families}\label{affine}
Let $W$ be an irreducible affine Coxeter group in its faithful Euclidean
action, and let $N\triangleleft W$ meet every maximal standard parabolic
trivially. The complete proof~\cite{translation} first shows that $N$ is
torsion-free: a finite-order isometry fixes a point, which can be moved
into the closed alcove and hence has a proper parabolic stabilizer.
For $g(x)=Ax+b\in N$, the commutator $gs_ig^{-1}s_i$ is the product
of reflections in $gM_i$ and $M_i$. Nonparallel mirrors give a nontrivial
finite-order element, so each wall-normal line is preserved by $A$.
Connectedness of the nonorthogonality graph of the wall normals forces
$A=I$ or $-I$. The latter makes $g$ an involution. Thus $N$ consists of
translations. This also covers the one-dimensional case directly.

Consequently source \emph{Corollary}~5.1 holds for the locally universal
condition explicitly imposed in its Section~5: the hypertope is an incidence
quotient by a normal translation subgroup. Maxwell's affine normal-subgroup
classification is not required by this proof. The conclusion is necessary,
not sufficient: arbitrary translation subgroups can fail~\eqref{pair}
or full flag validity. The affine $A_2$ example above also qualifies the
source's page~11 assertion that avoiding individual parabolics completes
quotient classification. Euclidean diagrams with projective residues, such
as the source's attributed Schulte/Hartley examples, remain outside this
locally universal kernel conclusion.

The displayed family classifications are retained with their exact domains.
Coxeter~\cite{CoxeterMaps} and Cacciari~\cite{Cacciari} give, up to duality,
\[
 \{4,4\}_{(s,0)},\ \{3,6\}_{(s,0)}\ (s\ge2),\qquad
 \{4,4\}_{(s,s)},\ \{3,6\}_{(s,s)}\ (s\ge1),
\]
and $(3,3,3)_{(s,0)},(3,3,3)_{(s,s)}$ with $s\ge2$ in both cases.
McMullen--Schulte~\cite{Toroids,Polytopes} supply the cubic families
$\{4,3^{n-2},4\}_{(s^k,0^{n-k})}$ for $s\ge2$, $k=1,2,n$, and
the $F_4$ families $\{3,3,4,3\}_{(s^k,0^{4-k})}$ for $s\ge2$, $k=1,2$.
These are attributed classification inputs, not new lattice exhaustions.

For the two nonstring rank-four diagrams the literal adapter~\cite{period}
proves existence and marked groups of all six displayed families. Put
$E=2\Z^3+\Z(1,1,1)$; $B_3^+$ is the signed permutation subgroup with an
even number of sign changes. Their groups are $(D/L)\rtimes B_3$ and
$(D/L)\rtimes B_3^+$, respectively:
\begin{center}\small
\begin{tabular}{llrr}\toprule
Period and domain &$L$&$|G_{\widetilde B_3}|$&$|G_{\widetilde A_3}|$\\\midrule
$(2s,0,0)$, $s\ge2$&$2s\Z^3$&$192s^3$&$96s^3$\\
$(s,s,0)$, $s\ge2$&$sD$&$48s^3$&$24s^3$\\
$(2s,2s,2s)$, $s\ge1$&$2sE$&$768s^3$&$384s^3$\\\bottomrule
\end{tabular}
\end{center}
Actual halving geometries, including the smallest allowed periods, supply
thinness, residual connectivity and full regularity. The claim that these
families exhaust the rank-four cases is retained from Ens~\cite{Ens}.
The source's announced all-rank geometric classification in its reference~21
is an announcement, not a proved result acquired here.

\section{Compact diagrams and finite existence}\label{hyperbolic}
The finite, affine and compact hyperbolic simplex diagram classifications
are retained from Humphreys~\cite{Humphreys} and
McMullen--Schulte~\cite{Polytopes}. The exact twenty-one printed rows of
source Table~2 are recorded in the clause ledger. Its ten affine rows are
$\widetilde A_{n-1}$ ($n\ge3$), $\widetilde D_{n-1}$ ($n\ge5$),
$\widetilde C_{n-1}$ ($n\ge3$), $\widetilde B_{n-1}$ ($n\ge4$),
$\widetilde E_6,\widetilde E_7,\widetilde E_8,\widetilde F_4,
\widetilde G_2,\widetilde A_1$. The last is rank two with an infinity bond.
The eleven hyperbolic rows comprise the strings $[p,q]$ with
$1/p+1/q<1/2$, and triangles $(k,l,m)$ with $1/k+1/l+1/m<1$;
the rank-four squares $(3,3,3,p)$, $p=4,5$, $(3,4,3,4)$,
$(3,5,3,p)$, $p=4,5$; strings $[p,3,5]$, $p=4,5$, and $[3,5,3]$;
the fork $[5,3_3]$; the rank-five pentagon $(3,3,3,3,4)$;
strings $[p,3,3,5]$, $p=3,4,5$; and the fork $[5,3,3_3]$.
All unspecified finite labels in this description are at least three.
Repeated parameter values give nine rank-four and five rank-five diagrams.

\begin{theorem}[Acquired finite-existence theorem]\label{exist}
If an infinite Coxeter system has finite rank at least three and every
maximal standard parabolic is finite, it has finite regular hypertope
quotients of unbounded chamber number with unchanged diagram and universal
spherical proper residues.
\end{theorem}
The proof~\cite{existence} avoids every nonidentity element in the finite
union $\bigcup_{i,j,k}P_iP_jP_k$. Two factors ensure the full intersection
property; three factors ensure that pairwise incident quotient objects
represented within one parabolic were incident before quotienting. Every
flag then lifts and each full residue incidence is preserved. In the
faithful geometric representation all coefficients lie in a fixed ring of
algebraic integers. A rational prime avoiding finitely many nonzero norms
separates this finite set by reduction. Separating arbitrarily many further
elements gives unbounded chamber numbers. This proves existence for all
retained compact diagrams, including nonstrings, but supplies no small-order
bound. Rank two is expressly excluded: a finite quotient of the infinite
dihedral group changes its infinity bond.

The earlier constructions remain attributed to their authors.
Corn--Singerman~\cite{Corn} provide the regular-hypermap triangle-rotation
presentation; Monson--Schulte~\cite{Monson} the modular $\{3,5,3\}$ examples;
Breda--Jones--Schulte~\cite{Breda} the $D_s^6\rtimes[3,5]$, $s\ge2$,
construction; Hartley--Leemans~\cite{Hartley} the $J_1\times J_1$ and
$J_1\times J_1\times L_2(19)$ examples. The star-polytope construction
of McMullen~\cite{Star}, also described by Coxeter~\cite{CoxeterPoly},
has the stated rank-four $120$ facets and vertex figures and group $H_4$
of order $14400$. Its deep-hole relations are retained as part of that
specific construction. The rank-five deep-hole object cited on source
page~13 is an infinite $5$-apeirotope; that page does not explicitly call
it finite. Theorem~\ref{exist} proves finite existence independently and
does not reproduce these particular groups or historical constructions.

\section{The complete finite-table disposition}\label{finite}
For a square $(a,b,c,d)$ the four cycle edges have those labels and the
diagonals label two. Strings $[a,b,c]$ have three consecutive bonds, and
$[5,3_3]$ is the rank-four fork with bonds $5,3,3$. This fixes the ordered
Coxeter presentations up to the source's permitted diagram relabelling.
Every regular hypertope with such a diagram gives a normal kernel in this
presentation of index equal to its number of chambers. No assumption of
universal local residues enters this necessary enumeration interface.

The eight-type core~\cite{core} completely checks all $120$ normal kernels
through its eight respective bounds and obtains $59$ positive marked
kernels. Its test checks all marked subgroup intersections and all sets of
objects incident with each base flag. Induction on flag size proves that
this tests \emph{all} actual flags, not only those already known to lie in
chambers. It therefore pays flag extension, thinness, residual connectivity,
and the full type-preserving group. The separate structural certificates
in~\cite{names} prove exactly the following twenty-two source group rows.
Here split actions are the actual certified actions; notation alone is not
a uniqueness or isomorphism proof.
\begin{center}\small
\begin{tabular}{llr}\toprule
Diagram & Certified source group & Order\\\midrule
$(3,3,3,4)$ & $2^4:A_5:2$ & 1920\\
$(3,3,3,4)$ & $(2^3:L_2(7))\times C_2$ & 2688\\
$(3,3,3,4)$ & $2^5:A_5:2$ & 3840\\
$(3,3,3,4)$ & $(4\mathbin{\cdot}(2^4:A_5)):2$ & 7680\\
$(3,3,3,4)$ & $2^6:A_5:2$ & 7680\\
$(3,3,3,5)$ & $(A_5\times A_5):2$ & 7200\\
$(3,4,3,4)$ & $2^{1+4}:S_3:S_3$ & 1152\\
$(3,4,3,4)$ & $S_6\times C_2$ & 1440\\
$(3,4,3,4)$ & $S_6\times S_3$ & 4320\\
$(3,4,3,4)$ & $(2^3:L_2(7))\times C_2^2$ & 5376\\
$(3,4,3,4)$ & $L_2(7)\times S_3^2$ & 6048\\
$(3,4,3,4)$ & $2^8:S_3:S_3$ & 9216\\
$[3,5,3]$ & $L_2(16):2$ & 8160\\
$[4,3,5]$ & $2^5:A_5:2$ & 3840\\
$[4,3,5]$ & $(2^4:A_5:2)\times C_2$ & 3840\\
$[4,3,5]$ & $(2^5:A_5:2)\times C_2$ & 7680\\
$[4,3,5]$ & $2^5:A_5:2:2$ & 7680\\
$[5,3,5]$ & $(A_5\times A_5):2$ & 7200\\
$[5,3,5]$ & $L_2(16):2$ & 8160\\
$[5,3_3]$ & $(2^5:A_5)\times C_2$ & 3840\\
$[5,3_3]$ & $2^6:A_5$ & 3840\\
$[5,3_3]$ & $(2^6:A_5)\times C_2$ & 7680\\
\bottomrule\end{tabular}\end{center}

The displayed nonsplit dot has an actual central/derived obstruction to a
complement. Split and direct factors have literal certificates, including
the affine $2^3:L_2(7)$ model, extraspecial $2^{1+4}$, and the outer
involution in $L_2(16):2$. Every row in this positive table has universal
spherical local groups.

Four printed rows are excluded even under the broad convention:
\begin{center}
\begin{tabular}{lll}\toprule
Diagram & Printed order & Complete exclusion\\\midrule
$(3,4,3,4)$&$336$&all three exact-order kernels fail IP\\
$(3,5,3,4)$&$2160$&the exact-order kernel fails IP\\
$(3,5,3,5)$&$3840$&all eleven exact-order kernels fail IP\\
$(3,5,3,5)$&$9720$&the exact-order kernel fails IP\\\bottomrule
\end{tabular}
\end{center}
The first two follow from~\cite{core}. The complete ninth-square search
through $3840$~\cite{square} returns $25$ kernels and $23$ markings after
the necessary rank-two divisibility filter. Twenty-one fail the intersection
property. The order-$960$ candidate satisfies it but has a nonextendible
three-type flag. The sole positive has order $1920$ and all four local
orders $60$. Its explicit permutation tuple and full incidence proof are
supplied; its global abstract structural name is not asserted.

For order $9720$, the exact-divisor method~\cite{divisor} preserves every
normal overgroup of a sought kernel: its index must divide $9720$.
The six returned indices are $1$, $2$, $60$, $120$, $4860$ and $9720$, each once.
Its sole exact-order quotient fails twelve marked intersections. An
independent integer closure exhibits the same nonidentity central
involution in all four maximal parabolics. This excludes that row without
repeating the timed-out unrestricted search through $9720$, and makes no
completeness claim at unrelated indices.

The complete bounded minima, including all source-equivalent markings, are:
\begin{center}\small
\begin{tabular}{lrrr}\toprule
Diagram & Complete bound & Broad minimum & Universal minimum\\\midrule
$(3,3,3,4)$ & 7680 & 168 & 1920\\
$(3,3,3,5)$ & 7200 & 2520 & 7200\\
$(3,4,3,4)$ & 9216 & 168 & 1152\\
$(3,5,3,4)$ & 2160 & 1080 & $>2160$\\
$(3,5,3,5)$ & 3840 & 1920 & $>3840$\\
$[3,5,3]$ & 8160 & 660 & 8160\\
$[4,3,5]$ & 7680 & 1920 & 3840\\
$[5,3,5]$ & 8160 & 3420 & 7200\\
$[5,3_3]$ & 7680 & 960 & 3840\\
\bottomrule\end{tabular}\end{center}

The broad minima are all attained. An inequality in the last column means
only exclusion through that bound; no attained universal minimum beyond it
is claimed. Source Table~3's assertion that each first row is smallest is
therefore replaced by this explicit two-convention table. For example the
product-size obstruction $48^2/6=384>336$ applies only to universal
$B_3$ residues; the full three-kernel test is what excludes $336$ without
that hypothesis. Kernel inventories retain all multiplicities and are
compared by exact ordered-generator isomorphisms, not compact permutation
byte equality or abstract group names.

These finite exhaustions retain the correctness and completeness contract
of GAP/LINS, its default normal-subgroup traversal, and its fixed tables of
characteristically simple groups, multiplier primes and socle overgroups.
The installed maximum index $10^7$ exceeds every bound used. The acquired
Coxeter seed and module-interface proofs~\cite{methods} and divisor proof
state the precise retained interfaces. Literal closures and coset readers
independently verify returned markings, but do not constitute a formal
verification of the normal-subgroup algorithm or its classification tables.

\section{Source disposition, premises and attribution}\label{closure}
The clause ledger contains $109$ entries: the original $83$ stable source
clauses, the $21$ explicit Table~2 rows, and five scope clauses. It routes
all numbered statements, both examples, all ten Table~1 rows, all affine
families, all twenty-six Table~3 rows and nine minimum assertions. Corrections
are local to their stated assertions: Theorem~3.2; local injectivity in the
broad reading of Theorem~3.4; $E_6$'s order; the overbroad Coxeter-element
half-power sentence; sufficiency of affine parabolic avoidance; and the
four finite-table rows and consequent minimum claims. The source's
rank-five small search remains an inconclusive historical search, not a
nonexistence theorem. No formal named conjecture is disposed of here.

The indispensable retained inputs are the classical Coxeter complex and
parabolic/chamber facts; faithful geometric representations and affine
alcove stabilizers; finite/affine/compact diagram classifications;
Maxwell's finite normal-subgroup alternative; the string correspondence;
the named toroidal classifications; the cited earlier special constructions;
and the exact finite software/classification-table contracts. Their roles
and source citations are explicit in the ledger. In particular, neither a
new all-rank affine lattice classification nor independent reconstructions
of every historical example are claimed. The source's abstract must be
read with its Section~5 review-and-announcement boundary.

The assembly checker verifies the exact acquired manuscripts, certificates,
reviews and their declared dependencies, the source-table joins, and the
native document. Small independent controls exercise the new direct-sum,
coset-kernel and $D_4$ joins and reject damaged ledgers. It does not rerun an
acquired group census or exceptional weight-orbit enumeration. Mathematical
validity of the new joins requires the separate reading of this proof;
hash identity is only an acquisition check.

The classifications and source questions belong to the cited authors.
Component authorship and distinct review roles are preserved in their
records. The present author previously contributed the eight-type finite
core, named-group certificates and Coxeter interfaces, and independently
reviewed the separate divisor, ninth-square and cubic applications.
The present synthesis and elementary joins require their own distinct
review. No claim of historical novelty, new formalization, new question
credit or measured compression follows from this assembly.

\begin{thebibliography}{99}
\bibitem{FLW} M.~E. Fernandes, D. Leemans and A. Ivi\'c Weiss,
\emph{An exploration of locally spherical regular hypertopes}. Fixed
sixteen-page revised manuscript, catalogue L033. Source reference numbers
are retained in the machine-readable bibliography ledger.
\bibitem{Maxwell} G. Maxwell, \emph{The normal subgroups of finite and affine
Coxeter groups}, Proc. London Math. Soc. 76 (1998), 359--382.
Theorem 0.1 retained through the fixed source.
\bibitem{Humphreys} J.~E. Humphreys, \emph{Reflection Groups and Coxeter
Groups}, Cambridge University Press, 1990, Section 6.9.
\bibitem{Polytopes} P. McMullen and E. Schulte, \emph{Abstract Regular
Polytopes}, Cambridge University Press, 2002, especially Sections 3C and 6B.
\bibitem{CoxeterMaps} H.~S.~M. Coxeter, \emph{Configurations and maps},
Rep. Math. Colloq. 8(2) (1948), 11--38.
\bibitem{Cacciari} L. Cacciari, \emph{On the automorphism group of a toroidal
hypermap}, Informatique th\'eorique et applications 26 (1992), 523--540.
\bibitem{Toroids} P. McMullen and E. Schulte, \emph{Higher toroidal regular
polytopes}, Adv. Math. 117 (1996), 17--51.
\bibitem{Ens} E. Ens, \emph{Rank 4 toroidal hypertopes}, Ars Math.
Contemp. 15 (2018), 67--79.
\bibitem{Corn} D. Corn and D. Singerman, \emph{Regular hypermaps},
European J. Combin. 9 (1988), 337--351.
\bibitem{Monson} B. Monson and E. Schulte, \emph{Modular reduction in abstract
polytopes}, Canad. Math. Bull. 52 (2009), 435--450.
\bibitem{Breda} A. Breda, G.~A. Jones and E. Schulte, \emph{Constructions
of chiral polytopes of small rank}, Canad. J. Math. 63 (2011), 1254--1283.
\bibitem{Hartley} M.~I. Hartley and D. Leemans, \emph{On locally spherical
polytopes of type $\{5,3,5\}$}, Discrete Math. 309 (2009), 247--254.
\bibitem{Star} P. McMullen, \emph{The groups of the regular star-polytopes},
Canad. J. Math. 50 (1998), 426--448, especially Sections 5--6.
\bibitem{CoxeterPoly} H.~S.~M. Coxeter, \emph{Regular Polytopes},
third edition, Dover, 1973.
\bibitem{residue} Human Compression project, \emph{Residue stabilizers and
distinguished generators}, complete reviewed L052 general manuscript.
\bibitem{kernel} Human Compression project, \emph{Kernels and intersections
in quotients of hypertopes}, L033 quotient-kernel component.
\bibitem{spherical} Human Compression project, \emph{Spherical hypertopes
from marked kernels and antipodal flags}, L033 spherical component, version 2.
\bibitem{translation} Human Compression project, \emph{Translation kernels
in affine reflection groups}, L033 translation component, version 2.
\bibitem{period} Human Compression project, \emph{The rank-four toroidal
families in cubic coordinates}, L033 period adapter.
\bibitem{existence} Human Compression project, \emph{Finite quotients
preserving spherical residues}, L033 finite-existence component, version 2.
\bibitem{core} Human Compression project, \emph{Eight bounded rank-four
hyperbolic quotient searches}, L033 finite core.
\bibitem{names} Human Compression project, \emph{Structural certificates
for twenty-two hyperbolic marked groups}, L033 named-group component.
\bibitem{methods} Human Compression project, complete reviewed Coxeter
seed and LINS composition interfaces, L033 method component.
\bibitem{square} Human Compression project, complete square $(3,5,3,5)$
quotients through order $3840$, L033 ninth-square component.
\bibitem{divisor} Human Compression project, exact-divisor Coxeter
quotient search and the order-$9720$ exclusion, L033 divisor component.
\bibitem{cubes} Human Compression project, infinite cubic example and
its exact dual-cell interface, L033 Example 2.5 component.
\end{thebibliography}
\end{document}
